\section{Fazit und Ausblick}
\label{sec:conclusion}

\subsection{Fazit}
\label{sec:fazit}

Ziel dieser Arbeit war die Entwicklung und Untersuchung eines
Architekturkonzepts für die kontrollierte KI-gestützte Transformation
heterogener Engineering-Bestandsinformation in SysML~v2. Ausgangspunkt war
die Frage, wie bestehende Engineering-Information durch probabilistische
Verfahren interpretiert und für die Modellbildung nutzbar gemacht werden
kann, während Herkunft, Unsicherheit und fachliche Entscheidungsbefugnis
nachvollziehbar erhalten bleiben.

Der wesentliche Erkenntnisgewinn liegt im Zusammenspiel der dafür
erforderlichen Architekturmechanismen. Eine kontrollierte Informationsbasis
unterscheidet Engineering-Quellen, Kontextquellen sowie den technischen
Verarbeitungs- und Orchestrierungskontext. Quellengebundene Evidenz verbindet
fachliche Aussagen mit ihrer Informationsgrundlage. Mehrere
Persona-Perspektiven interpretieren gemeinsame fachliche
Referenzgegenstände und stellen Übereinstimmung, Varianz und gemeinsame
Mehrdeutigkeit als zusätzliche Bewertungsinformation bereit. Die daraus
entstehende Engineering-Information wird anschließend durch explizite
menschliche Entscheidungen autorisiert und kontrolliert in die Modellbildung
überführt.

Dabei hat sich insbesondere die Trennung der einzelnen
Entscheidungsgegenstände als wesentlich erwiesen. Die Human Engineering
Authority autorisiert die fachliche Verwendung einer Information, während die
Model Authority über ihre strukturelle Einordnung und Repräsentation im
Zielmodell entscheidet. Fachliche Bedeutung, Modellrepräsentation und
zielmodellspezifische Formulierung bleiben dadurch unterscheidbare
Verarbeitungs- und Entscheidungsebenen. Nach den erforderlichen
Autorisierungen kann die weitere Überführung regelgebunden über einen
kontrollierten internen Modellzustand bis zur SysML-v2-Repräsentation
erfolgen. Technische Validierung, menschliche Modellprüfung und Freigabe
bilden daran anschließend getrennte Prüf- und Entscheidungsstufen.

Ein weiteres zentrales Ergebnis ist die durchgängige Provenienz der
verarbeiteten Engineering-Information. Quellen, maschinelle
Interpretationen, menschliche Autorisierungsentscheidungen und
Modellzustände bleiben als eigenständige und miteinander verbundene Zustände
erhalten. Auf dieser Grundlage kann die Entstehung eines resultierenden
Modellzustands von seiner Informationsgrundlage bis zur
SysML-v2-Repräsentation und in umgekehrter Richtung zurück zu den zugrunde
liegenden Informations- und Entscheidungszuständen nachvollzogen werden.
Provenienz und Nachverfolgbarkeit bilden damit keine nachträgliche
Dokumentation des Ergebnisses, sondern eine querschnittliche Eigenschaft des
gesamten Transformationsprozesses.

Die prototypische Untersuchung hat zugleich die Bedeutung enger
Kontrollmechanismen für den Einsatz generativer KI im Engineering
verdeutlicht. Probabilistische Verfahren können heterogene Information
erschließen, unterschiedliche Interpretationen erzeugen und damit
Engineering-Aufgaben wirkungsvoll unterstützen. Ihre Ergebnisse erhalten
dadurch jedoch keine eigenständige fachliche Entscheidungsbefugnis. Der
Mensch wird deshalb nicht lediglich als abschließender Prüfer eines bereits
erzeugten Modells eingebunden, sondern entscheidet an den Übergängen, an
denen sich die fachliche Bedeutung oder die modellbezogene Wirkung einer
Information verändert. Die kontrollierte KI-Unterstützung entsteht damit aus
dem Zusammenspiel probabilistischer Verarbeitung, expliziter menschlicher
Autorisierung und regelgebundener Weiterverarbeitung.

Die formative End-to-End-Erprobung war für diese Erkenntnis von besonderer
Bedeutung. Mehrere zunächst zusammengefasste Informations- und
Entscheidungsfunktionen wurden erst durch die praktische Ausführung des
Transformationsprozesses als eigenständige Architekturgrenzen sichtbar und
anschließend präzisiert. Die Verarbeitung mehrerer Engineering-Quellen
ergänzte diese Untersuchung um den Übergang von quellenbezogen autorisierten
Informationszuständen zur gemeinsamen projektweiten Modellbildung. In der
abschließenden Demonstration konnte der zusammenhängende Verarbeitungspfad
von den registrierten Quellen über KI-gestützte Interpretation und
menschliche Entscheidungen bis zum technisch validierten und menschlich
freigegebenen SysML-v2-Artefakt ausgeführt und anhand der persistierten
Zustände anschließend rekonstruiert werden. Das entwickelte
Architekturkonzept ist damit nicht ausschließlich Ergebnis einer
konzeptionellen Herleitung, sondern wurde durch seine prototypische
Operationalisierung und Untersuchung schrittweise präzisiert.

Eine weitergehende Erkenntnis betrifft den Aufwand, der einem sinnvollen
KI-Einsatz vorausgeht. Bevor probabilistische Verfahren zielgerichtet in
einem Engineering-Prozess eingesetzt werden können, müssen unter anderem die
Informationsgrundlage, die zulässigen Informationsrollen,
Entscheidungsbefugnisse, Modellierungsregeln und Prüfmechanismen festgelegt
werden. Ein wesentlicher Teil der Gestaltungsarbeit liegt damit vor der
eigentlichen KI-gestützten Verarbeitung. Die entwickelte Architektur bündelt
mehrere dieser wiederkehrenden Entscheidungen in einem zusammenhängenden
Konzept und stellt damit einen bereits ausgearbeiteten Ausgangspunkt für
weitere Anwendungen bereit.

Damit schließt sich der Bogen zu dem in
Abschnitt~\ref{sec:problemstellung} dargestellten Zielbild einer
KI-unterstützten Brownfield-Transformation. Der dort skizzierte
Transformationspfad beginnt zunächst mit zusätzlichem Gestaltungs- und
Integrationsaufwand, bevor wiederverwendbare Verarbeitungsschritte ihre
Wirkung entfalten können. Arbeiten wie die vorliegende können dazu beitragen,
einen Teil dieses konzeptionellen Initialaufwands in wiederverwendbare
Architekturmechanismen zu überführen. Werden solche Mechanismen weiter
konsolidiert, empirisch untersucht und für unterschiedliche
Engineering-Aufgaben verfügbar gemacht, entsteht die Perspektive, dass
zukünftige KI-gestützte Anwendungen auf einer zunehmend etablierten
architektonischen Grundlage aufbauen können.

Der wissenschaftliche Beitrag dieser Arbeit liegt damit nicht allein in der
Erzeugung einer SysML-v2-Repräsentation. Er besteht in der Entwicklung und
prototypischen Untersuchung einer Informationsverarbeitungsarchitektur, die
zeigt, wie heterogene Engineering-Bestandsinformation, probabilistische
Interpretation, menschliche Engineering-Entscheidungen, Provenienz,
Nachverfolgbarkeit und formale Modellbildung zu einem kontrollierten
Transformationsprozess verbunden werden können.

Mehrere dieser Mechanismen erfüllen dabei Funktionen, die nicht unmittelbar
an die SysML-v2-Zielrepräsentation gebunden sind. Insbesondere
Quellenverarbeitung, Provenienz, persistierte Verarbeitungszustände,
menschliche Autorisierung und die Trennung zwischen maschinellen Vorschlägen
und autorisierten Zuständen bilden einen wiederverwendbaren strukturellen
Kern. Daraus ergibt sich die weiterführende Frage, inwieweit dieser Kern auch
für andere KI-gestützte Engineering-Aufgaben genutzt werden kann. Diese
Perspektive bildet den Ausgangspunkt des folgenden Ausblicks.

\subsection{Ausblick}
\label{sec:ausblick}

Aus den Ergebnissen dieser Arbeit ergeben sich drei aufeinander aufbauende
Entwicklungsperspektiven. Die unmittelbarste betrifft die Weiterentwicklung
des untersuchten Prototyps für einen produktiven Einsatz im industriellen
Umfeld. Darüber hinaus stellt sich die Frage, welche Teile der entwickelten
Architektur auf andere Engineering-Aufgaben übertragen werden können.
Langfristig eröffnet die persistierte menschliche Interaktion schließlich die
Möglichkeit, die KI-gestützte Verarbeitung kontrolliert aus früheren
Entscheidungen lernen zu lassen.

\paragraph{Weiterentwicklung zu einem produktiv einsetzbaren System}

Der zunächst naheliegendste Entwicklungspfad besteht darin, den
prototypischen Stand innerhalb des untersuchten SysML-v2-Anwendungsfalls
weiterzuführen. Ausgangspunkt hierfür können die während Entwicklung und
Evaluation dokumentierten Befunde bilden. Eine Triage der im Anhang
festgehaltenen offenen Punkte könnte zunächst unterscheiden, welche Befunde
für einen produktiven Einsatz funktional, technisch oder hinsichtlich der
Bedienbarkeit relevant sind und in welcher Reihenfolge sie adressiert werden
sollten.

Darauf aufbauend wären insbesondere diejenigen Aspekte weiterzuentwickeln,
die innerhalb des prototypischen Einsatzes nur in begrenztem Umfang
untersucht wurden. Dazu gehören die Verarbeitung umfangreicherer und weniger
stark aufbereiteter Informationsbestände, zusätzliche
SysML-v2-Modellierungskonstrukte, die technische und fachliche Behandlung
weiterführender Zustandsänderungen sowie die Skalierung der menschlichen
Prüf- und Autorisierungsschritte. Ergänzend könnten Softwarequalität,
Laufzeit, Ressourcenbedarf und Robustheit des Gesamtsystems systematischer
untersucht werden.

Ein wesentlicher nächster Schritt wäre außerdem die Einbettung des Systems in
einen konkreten Engineering-Workflow eines Unternehmens. Dabei könnten die
im Prototyp verwendeten fachlichen Referenzen, Beispielmodelle,
Persona-Perspektiven, Regeln und Modellierungsrahmen an den jeweiligen
Unternehmenskontext angepasst werden. Der modulare Aufbau der Architektur
ermöglicht eine solche Konfiguration, ohne den gesamten
Transformationsprozess neu entwerfen zu müssen. In einer solchen
Anwendungsumgebung könnte anschließend untersucht werden, wie sich das System
mit realen Engineering-Quellen, bestehenden Modellierungsrichtlinien und den
tatsächlichen Entscheidungsprozessen der beteiligten Systems Engineers
verbindet.

Eine weiterführende Validierung mit unabhängigen Anwendern könnte dabei
insbesondere untersuchen, ob die bereitgestellten
Entscheidungsinformationen verständlich sind, wie hoch der tatsächliche
Prüfaufwand ausfällt und wie sich die vorgesehenen Autorisierungsgrenzen in
bestehende Arbeitsabläufe integrieren lassen. Auf diese Weise könnte der
prototypisch untersuchte Transformationsprozess schrittweise zu einem
industriell einsetzbaren Werkzeug weiterentwickelt werden.

\paragraph{Übertragung auf weitere Engineering-Aufgaben}

Eine zweite Entwicklungsperspektive ergibt sich aus der Beobachtung, dass
mehrere zentrale Architekturmechanismen nicht unmittelbar an die
SysML-v2-Zielrepräsentation gebunden sind. Quellenverarbeitung,
Provenienzbeziehungen, stabile Identitäten und Verarbeitungszustände,
probabilistische Interpretation, menschliche Autorisierung sowie die
Trennung zwischen maschinellen Vorschlägen und autorisierten Zuständen werden
auch für andere dokumentenzentrierte Engineering-Aufgaben benötigt.

Diese Perspektive wird im Abgabestand des Prototyps unter dem Arbeitsbegriff
\emph{Reusable Engineering Document Processing Architecture} (REDPA)
konkretisiert. Hierfür liegen bereits Adaptionsbeschreibungen für drei
alternative Anwendungsszenarien vor. Die Dokumente unterscheiden jeweils
zwischen generischen Bestandteilen, die weiterverwendet werden können, und
fach- beziehungsweise zielspezifischen Bestandteilen, die konfiguriert,
ersetzt oder ergänzt werden müssten. Zu den vorgesehenen wiederverwendbaren
Mechanismen gehören insbesondere Quellenverarbeitung, Provenienz,
Identitäts- und Fingerprintmechanismen, persistierte Verarbeitungszustände,
die technische Abstraktion der Sprachmodelle sowie menschliche Prüf- und
Freigabemechanismen. :contentReference[oaicite:1]{index=1}
:contentReference[oaicite:2]{index=2}
:contentReference[oaicite:3]{index=3}

Ein dokumentierter Adaptionspfad beschreibt einen
\emph{Requirements Review Assistant}, der Requirements quellengebunden
analysiert, mögliche Qualitäts- und Konsistenzbefunde erzeugt und diese einer
menschlichen Entscheidung zuführt. Ein weiterer Pfad beschreibt einen
\emph{Technical Documentation SOP Tracker}, der Verpflichtungen aus
Verfahrens- und SOP-Dokumenten mit vorhandener technischer Dokumentation
verknüpft und daraus einen geprüften Nachverfolgungszustand ableitet. Der
dritte Pfad überträgt die Architektur auf einen
\emph{QM/RA Documentation Reviewer}, der Dokumentation gegen konfigurierbare
Qualitäts-, regulatorische und interne Regelwerke prüft und die erzeugten
Befunde quellengebunden für eine menschliche Bewertung bereitstellt.
:contentReference[oaicite:4]{index=4}
:contentReference[oaicite:5]{index=5}
:contentReference[oaicite:6]{index=6}

Diese Adaptionsbeschreibungen konkretisieren die Hypothese eines
wiederverwendbaren Verarbeitungskerns, stellen jedoch noch keine Untersuchung
seiner Übertragbarkeit dar. Eine weiterführende Arbeit könnte deshalb einen
oder mehrere dieser Pfade realisieren und systematisch untersuchen, welche
Architekturelemente tatsächlich unverändert übernommen werden können und
welche aufgrund der jeweiligen Engineering-Aufgabe angepasst werden müssen.
Auf diese Weise könnte eine REDPA schrittweise aus den in dieser Arbeit
entwickelten Mechanismen abstrahiert und empirisch konsolidiert werden.

Eine solche Wiederverwendung würde zugleich den in
Abschnitt~\ref{sec:problemstellung} beschriebenen Gedanken aufgreifen, dass
vor einem zielgerichteten KI-Einsatz zunächst ein geeigneter
Verarbeitungsrahmen geschaffen werden muss. Mit einer zunehmend
konsolidierten Grundarchitektur könnte sich dieser Gestaltungsaufwand stärker
auf die jeweilige fachliche Aufgabe, ihre Regeln und ihre
Entscheidungsgegenstände konzentrieren.

\paragraph{Lernen aus menschlichen Engineering-Entscheidungen}

Eine dritte, weiterführende Perspektive ergibt sich aus den im
Transformationsprozess persistierten menschlichen Entscheidungen. Die
Human-in-the-Loop-Mechanismen erfüllen zunächst die Funktion, maschinell
erzeugte Interpretationen und Modellentscheidungen kontrolliert zu
autorisieren. Gleichzeitig entsteht durch jede dieser Interaktionen ein
strukturierter Zusammenhang zwischen der zugrunde liegenden
Engineering-Information, einem maschinellen Vorschlag und seiner
menschlichen Bewertung oder Korrektur.

Mit zunehmender Nutzung entsteht dadurch ein wachsender Bestand an
fachspezifisch gelabelten Daten. Neben angenommenen und zurückgewiesenen
Vorschlägen sind insbesondere menschliche Änderungen sowie bewusst
unaufgelöst gebliebene Fälle relevant. Vergleichbare Ansätze des Lernens aus
menschlichen Präferenzen beziehungsweise menschlichem Feedback zeigen
grundsätzlich, dass solche Rückmeldungen zur Anpassung des Verhaltens
lernender Systeme genutzt werden können
\cite{Christiano2017,Ouyang2022}.

Für eine zukünftige Erweiterung des hier entwickelten Systems könnte daraus
ein kontrollierter Feedbackmechanismus entstehen. Als eine mögliche
Bewertungsgröße lässt sich dabei das \emph{Akzeptanzniveau} definieren. Es
beschreibt, in welchem Umfang maschinell erzeugte Vorschläge innerhalb eines
definierten Aufgabenbereichs durch die Human Engineering Authority ohne oder
nur mit geringer fachlicher Änderung übernommen werden. Ergänzend liefert die
bereits im Architekturkonzept betrachtete Varianz zwischen den
Persona-Perspektiven Information darüber, wie einheitlich oder unterschiedlich
ein fachlicher Gegenstand interpretiert wird.

Beide Größen könnten gemeinsam mit den tatsächlich vorgenommenen
Korrekturen als Lernsignal verwendet werden. Ziel wäre dabei nicht eine
pauschale Minimierung der Varianz. Unterschiedliche Perspektiven können
berechtigte Mehrdeutigkeit oder fachliche Unsicherheit sichtbar machen und
sollen in solchen Fällen erhalten bleiben. Ein lernendes System sollte
vielmehr erreichen, dass vermeidbare Abweichungen zurückgehen und
maschinelle Vorschläge häufiger mit der späteren fachlichen Entscheidung
übereinstimmen. Ein steigendes Akzeptanzniveau bei gleichzeitig angemessener
Abbildung tatsächlich vorhandener Varianz könnte hierfür eine mögliche
Bewertungsgrundlage darstellen.

Wiederkehrende menschliche Korrekturen könnten anschließend beispielsweise
zur Anpassung von Regeln, Beispielen, Interpretationsanweisungen oder
spezialisierten Modellen genutzt werden. Die bereits entwickelte Provenienz
ermöglicht es dabei, ein Lernsignal mit seinem jeweiligen fachlichen
Ausgangskontext und der tatsächlich getroffenen menschlichen Entscheidung zu
verbinden. Neue oder angepasste Verarbeitungsstände könnten versioniert,
zunächst außerhalb der aktiven Verarbeitung untersucht und erst nach einer
expliziten Freigabe übernommen werden.

Mit zunehmender Qualität der maschinellen Interpretation könnte sich dadurch
langfristig auch die Rolle des Menschen verändern. Statt jeden
Informationszustand in gleicher Intensität zu prüfen, könnte sich die
menschliche Aufmerksamkeit stärker auf Fälle mit hoher Unsicherheit,
relevanter Varianz, widersprüchlicher Informationslage oder geringer
bisheriger Akzeptanz konzentrieren. Ob und unter welchen Bedingungen einzelne
Entscheidungen schließlich weiter automatisiert werden können, wäre dabei
anhand der über die Nutzung entstehenden Daten empirisch zu untersuchen.

Die drei Entwicklungspfade adressieren damit unterschiedliche
Reifegrade derselben Grundidee. Kurzfristig kann der bestehende Prototyp für
einen produktiven Einsatz im SysML-v2-Kontext erweitert und in reale
Engineering-Workflows integriert werden. Darauf aufbauend kann untersucht
werden, welche Architekturmechanismen sich über REDPA auf weitere
Engineering-Aufgaben übertragen lassen. Langfristig eröffnet die dabei
entstehende strukturierte menschliche Rückmeldung die Möglichkeit eines
kontrolliert lernenden Systems, dessen zunehmende Automatisierung weiterhin
auf nachvollziehbaren Informationsgrundlagen und explizit überprüften
Engineering-Entscheidungen aufbaut.
